% Generated draft: requires mathematical and source review.
\documentclass[11pt]{article}
\usepackage[margin=25mm]{geometry}
\usepackage{fontspec}
\setmainfont{Latin Modern Roman}
\usepackage{amsmath,amssymb,mathtools,unicode-math}
\setmathfont{Latin Modern Math}
\usepackage{enumitem,needspace,xurl,hyperref}
\setlength{\parindent}{0pt}
\setlength{\parskip}{5pt}
\setlength{\emergencystretch}{3em}
\newcommand{\N}{\mathbb{N}}
\newcommand{\Z}{\mathbb{Z}}
\newcommand{\Q}{\mathbb{Q}}
\newcommand{\R}{\mathbb{R}}
\newcommand{\C}{\mathbb{C}}
\newcommand{\F}{\mathbb{F}}
\newcommand{\A}{\mathbb{A}}
\newcommand{\PP}{\mathbb{P}}
\newcommand{\E}{\mathbb{E}}
\newcommand{\eps}{\varepsilon}
\newcommand{\dd}{\,\mathrm{d}}
\begin{document}
\section*{UnsolvedMath ID: 800701}
\textbf{Source problem code:} catalogue entry
\subsection{Definitions and mathematical statement}
Let \(F\) be a field. A central simple \(F\)-algebra is a finite-dimensional associative \(F\)-algebra with identity, center \(F\), and no nonzero proper two-sided ideals. Its degree is \(\deg(A)=\sqrt{\dim_F A}\). The Brauer group \(\operatorname{Br}(F)\) consists of equivalence classes of central simple \(F\)-algebras, where two algebras are equivalent if suitable matrix algebras over them are isomorphic; its operation is induced by tensor product over \(F\). Write \([A]\) for the class of \(A\), and let \(\operatorname{per}(A)\) be the order of \([A]\) in \(\operatorname{Br}(F)\). Thus the source's word ``similar'' means equality of Brauer classes.

A cyclic algebra of degree \(m\) is a cyclic crossed product \((L/F,\sigma,a)\), where \(L/F\) is a cyclic Galois field extension of degree \(m\), \(\sigma\) generates its Galois group, and \(a\in F^\times\). It is generated by \(L\) and an element \(u\) satisfying \(u^m=a\) and \(ux=\sigma(x)u\) for \(x\in L\).

Let \(m,n\) be positive integers with the same prime divisors and with \(m\mid n\). Interpreting the requested number as a uniform bound on the number of factors, define
\[
\ell_F(m,n)=\min\left\{k\in\Z_{\geq0}:\begin{array}{l}
\text{for every central simple }F\text{-algebra }A\text{ with}\ \deg(A)=n\ \text{and}\ \operatorname{per}(A)\mid m,\\
\text{there exist }0\leq t\leq k\text{ cyclic }F\text{-algebras }C_1,\ldots,C_t\\
\text{of degree }m\text{ such that }[A]=[C_1\otimes_F\cdots\otimes_F C_t]
\end{array}\right\}.
\]
The empty tensor product represents the zero Brauer class. Set \(\ell_F(m,n)=\infty\) if no finite uniform bound exists. The source literally uses a tensor product of \(\ell_F(m,n)\) factors; whether it requires exactly that many factors, and which split cyclic factors are permitted for padding, is not specified.

The source assumes that \(\operatorname{Br}_p[m](F)\) is generated by cyclic algebras. The notation \(\operatorname{Br}_p[m](F)\), including the role of its subscript, is not defined in the supplied record, so this hypothesis is retained without identifying it with a particular subgroup of \(\operatorname{Br}(F)\).

\textbf{Problem.} Let \(p\) be a prime and let \(r\leq s\) be integer exponents for which \(p^r,p^s\) are positive integers. Suppose that \(\operatorname{Br}_p[p^r](F)\) is generated by cyclic algebras, in the source's unresolved notation. Determine
\[
\ell_F(p^r,p^s).
\]
The source does not explicitly specify whether the boundary case \(r=0\) is included.

As background, the supplied source states that if \(F\) contains the \(m\)-th roots of unity, the generic division algebra and the Merkurjev--Suslin theorem show that \(\ell_F(m,n)\) is finite. It further states that in this case one may reduce to \(m\) a prime power by [Tig84, Th. 2.2], and calls the invariant the \emph{symbol length}. These are supplied contextual assertions, not independently verified here. The roots-of-unity condition is not explicitly imposed in the final problem, and no characteristic restriction on \(F\) is supplied.
\subsection{Short English statement}
For a field \(F\), determine the least uniform number of cyclic algebras of degree \(p^r\) needed to represent every Brauer class having a representative of degree \(p^s\) and period dividing \(p^r\), under the source's cyclic-generation hypothesis. The supplied notation for that hypothesis and some boundary conventions require clarification.
\subsection{Sources}
\begin{itemize}
\item Local UnsolvedMath registry, permanent ID \(800701\), exact source problem code \texttt{catalogue entry}; title: \emph{Open Problems on Central Simple Algebras --- Problem 11}. The registry attributes the problem to Auel, Brussel, Garibaldi, and Vishne, with proposed year 2010, and labels it open. This is a supplied status label, not an independently verified assessment.
\item Auel, Brussel, Garibaldi, and Vishne, \emph{Open Problems on Central Simple Algebras} (2010), Problem 11. Supplied source URL: \url{https://arxiv.org/abs/1006.3304}. The record reports access on 2026-10-16 and extraction by \texttt{arxiv\_latex}; these are supplied provenance statements, not a claim of access during this preparation.
\item The supplied excerpt invokes the Merkurjev--Suslin theorem for its finiteness assertion, without a separate bibliographic reference.
\item The supplied excerpt cites [Tig84, Th. 2.2] for reduction to prime powers. No full bibliographic entry or link identifying [Tig84] is supplied.
\end{itemize}
\end{document}
